
\Problem{Tesoros Inaccesibles}

\Statement

Un árbol binario de \verb|char|s representa el mapa de una cueva. Cada nodo indica una de tres cosas:

\begin{itemize}
\item \verb|'.'|: hay un pasillo que se puede cruzar.
\item \verb|'X'|: hay un pasillo que \textbf{no} se puede cruzar.
\item \verb|'*'|: hay un tesoro (y se puede cruzar).
\end{itemize}

\medskip

A la cueva se entra por la raíz del árbol. Un tesoro es \textbf{accesible} si se puede llegar a él desde la raíz sin pasar por ningún nodo con \verb|'X'|. Un tesoro es \textbf{inaccesible} si algún nodo del camino desde la raíz hasta el tesoro contiene una \verb|'X'|.

\medskip

Fíjate en que la raíz del árbol podría ser una \verb|'X'|, y en ese caso todos los tesoros del árbol serían inaccesibles.

\medskip

\textbf{Apartado 1.} Haz una función

\begin{verbatim}
int compta_tresors_accessibles(BinTree<char> t);
\end{verbatim}

que cuente el número de tesoros accesibles del árbol \verb|t|.

\medskip

\textbf{Apartado 2.} Haz una función

\begin{verbatim}
int compta_tresors_inaccessibles(BinTree<char> t);
\end{verbatim}

que cuente el número de tesoros inaccesibles del árbol \verb|t|.

\medskip

Nótese que un tesoro inaccesible es un nodo \verb|'*'| al que no se puede llegar desde la raíz porque algún nodo del camino es una \verb|'X'|. Concretamente, si un nodo contiene \verb|'X'|, todo el subárbol que cuelga de este nodo es inaccesible, y los tesoros que haya ahí se deben contar.

\Observation

En los archivos públicos (icono del gatito) encontrarás: \verb|main.cc|, \verb|tresors.cc| (donde debes implementar las funciones), \verb|bintree.hh|, \verb|bintree-io.hh|, \verb|assert.hh|, un \verb|Makefile| y el directorio \verb|.vscode|.

\medskip

En el archivo \verb|tresors.cc| hay dos funciones de esqueleto comentadas que debes descomentar para poder compilar y probar el programa. Puedes implementarlas por separado: hay juegos de prueba específicos para cada una, de manera que puedes validar un apartado sin tener terminado el otro.

\medskip

Debes implementar las dos funciones en el archivo \verb|tresors.cc| y enviar solo este archivo.

\Input

La primera línea de la entrada contiene \verb|apartat1| o \verb|apartat2|, indicando qué apartado se debe ejecutar. A continuación se leen árboles binarios de \verb|char|s en formato visual, uno tras otro.

\Output

Para cada árbol, se escribe el resultado de la función correspondiente.

\SampleTwoCol
